\setcounter{footnote}{97}

\printedpage{124}
% head: A. Grothendieck : Dérivateurs, Ch. VI, Édition des Universités de Montpellier II et Paris VII
% foot: Version 9 novembre 2001

$W_S^{\mathrm{d}}$, $W_S^{\mathrm{d}}$, and the ``forget $X$'' functor from $\underline{\mathrm{Liss}}_W X$ (even from $\mathrm{Cat}/X$) to $\mathrm{Cat}/S$, which is compatible with $W_X^{\mathrm{d}}$ and $W_S^{\mathrm{d}}$. We thus get a functor % [sic] W_S^d, W_S^d -> apparently W_S^d, W_X^d (or W_X^d, W_S^d), the same symbol is printed twice
\[
\left\{\begin{aligned}
&\mathrm{mult}_{X/S} : \mathrm{HOT}_W(S) \longrightarrow \mathrm{HOT}_W(S)\\
&\mathrm{mult}_{X/S}(\mathrm{hot}_{S,W}(Y)) \overset{\text{def}}{=} \mathrm{hot}_{S,W}(X\times_S Y) \qquad \text{if $Y$ \textit{smooth} over $S$}
\end{aligned}\right.
\]
(\footnote{This formula is also valid if $Y$ is in $\underline{\mathrm{Fib}}_W\, S$, cf. $(*)$, p. 202.}), since $\mathrm{HOT}_W(S)$ can be obtained as the localized category of both $\underline{\mathrm{Liss}}_W S$ and $\mathrm{Cat}/S$, for $W_S^{\mathrm{d}}$. When $X$ is itself $W$-smooth over $S$, this functor is none other than
\[
\left\{\begin{aligned}
&\mathrm{mult}_\xi : \eta \longmapsto \xi\times_S \eta,\\
&\text{where }\ \xi = \mathrm{hot}_{S,W}(X).
\end{aligned}\right.
\]

But this is not so if $X$ is not smooth, i.e. we must distinguish in this case, even for $Y$ $W$-smooth over $S$,
\[
\mathrm{hot}_{S,W}(X\times_S Y) \qquad\text{and}\qquad \mathrm{hot}_{S,W}(X)\times_S \mathrm{hot}_{S,W}(Y).
\]

The second term can be made explicit as $\mathrm{hot}_{S,W}(\Phi(X)\times_S Y)$ which receives $\mathrm{hot}_{S,W}(X\times_S Y)$, i.e.
\msp{page 201}
we have a canonical homomorphism
\[
\underbrace{\mathrm{hot}_{S,W}(X\times_S Y)}_{=\,\mathrm{mult}_{X/S}(Y)} \longrightarrow \underbrace{\mathrm{hot}_{S,W}(X)\times_S \mathrm{hot}_{S,W}(Y)}_{=\,\mathrm{hot}_{S,W}(\Phi(X)\times_S Y)},
\]
induced by
\[
X\times_S Y \longrightarrow \Phi(X)\times_S Y,
\]
but this last arrow has no reason to be in $W_S^{\mathrm{d}}$, and thus does not necessarily induce an isomorphism in $\mathrm{HOT}_W(S)$. (Counterexample page 154.)

We must look separately at the case \textit{where $X$ is in} $\underline{\mathrm{Fib}}_W S$, I feel like saying that in this case we have
\[
\mathrm{hot}_{S,W}(X\times_S Y) \xrightarrow{\ \sim\ } \underbrace{\mathrm{hot}_{S,W}(X)\times_S \mathrm{hot}_{S,W}(Y)}_{=\,\mathrm{hot}_{S,W}(\Phi(X)\times_S\, \Phi(Y))} \tag*{$(*)$} % doubt resolved in DER_VI-c_corrections.csv; input note: under the brace "hot" is printed larger than "(Phi(X)x_S Phi(Y))", which stands at subscript level -- argument of hot_{S,W} or part of its subscript
\]
(\footnote{This formula is established only by assuming either $Y$ $W$-smooth (in addition to $X$ $W$-fibered over $S$), or $X$ $W$-proper, e.g. $X$ $W$-perfect.}) for \textit{every} $Y$ in $\mathrm{Cat}/S$ (not necessarily $W$-smooth), a fortiori that in this case
\[
\left\{\begin{aligned}
&\mathrm{mult}_{X/S} = \mathrm{mult}_\xi \qquad \text{in } \mathrm{HOT}_W(S)\\
&\text{where } \xi = \mathrm{hot}_{S/W}(X/S) \text{ (if $X$ in $\underline{\mathrm{Fib}}_W S$).} % [sic] hot_{S/W} -> hot_{S,W}, the notation used everywhere else
\end{aligned}\right. \tag*{$(**)$}
\]

\printedpage{125}
% head: A. Grothendieck : Dérivateurs, Ch. VI, Édition des Universités de Montpellier II et Paris VII
% foot: Version 9 novembre 2001

In the case where $Y$ is smooth over $S$ (which suffices to establish $(**)$), it suffices to recall that $X\longrightarrow \Phi(X)$ is in $W_S^{\mathrm{f}}$ (follows from the fact that $X$ is $W$-fibered over $S$ (\footnote{even $W$-presmooth over $S$ suffices}), hence $X_s \longrightarrow s\backslash X$ is in $W$, $\forall\ s\in S$), hence $\underbrace{X\times_S Y}_{X_Y} \longrightarrow \underbrace{\Phi(X)\times_S Y}_{\Phi(X)_Y}$ is also in $W_Y^{\mathrm{f}}$, hence (since $X_Y$, $\Phi(X)_Y$ are $W$-fibered over $Y$) in $W_Y^{\mathrm{d}}$, hence also in $W_S^{\mathrm{d}}$. In fact, even without assuming $Y$ smooth, this shows that we have
\msp{page 202}
\[
\mathrm{hot}_{S,W}(X\times_S Y) \xrightarrow{\ \sim\ } \mathrm{mult}_{Y/S}(\mathrm{hot}_{S,W}(X)) \qquad \text{if $X$ is in $\underline{\mathrm{Fib}}_W S$.} \tag*{$(*)$}
\]

To have $(*)$ for arbitrary $Y$, we can by the foregoing assume $X$ smooth (replacing it by $\Phi(X)$ if need be), hence $X$ in $\underline{\mathrm{Liss}}_{0,W}S$. We must see that
\[
X\times_S Y \longrightarrow X\times_S \Phi(Y)
\]
is still in $W_S^{\mathrm{d}}$. Now if $X$ is \textit{$W$-proper} over $S$, the base change $X\longrightarrow S$ is compatible with $W_S^{\mathrm{d}}$, $W_X^{\mathrm{d}}$ (on all of $\mathrm{Cat}/S$), and we win. (This must be false without an assumption of $W$-properness on $X$, \textit{or} of $W$-smoothness on $Y$ \dots)

When we have
\[
f : S' \longrightarrow S,
\]
then
\[
f^* : \mathrm{HOT}_W(S) \longrightarrow \mathrm{HOT}_W(S')
\]
is compatible with the product structures.

\subsection*{\protect\cercle{8} The functors $f_!$: adjunction theorem, properties of the derivator $\mathrm{HOT}_W$.}

Let
\[
f : S' \longrightarrow S
\]
be in Cat, then the induced functor
\[
f_! : \mathrm{Cat}/S' \longrightarrow \mathrm{Cat}/S
\]
is compatible with the localizers $(W_{S'}^{\mathrm{d}}, W_S^{\mathrm{d}})$, or $(W_{S'}^{\mathrm{g}}, W_S^{\mathrm{g}})$, or $(W_{S'}, W_S)$. We thus get an induced functor for the localized categories, and \texttt{[a]} commutative diagram
\msp{page 203}
\[
\begin{tikzcd}[column sep=small]
\mathrm{HOT}_W(S') \arrow[rrd, "\rho'"] \arrow[dd, "f_!"'] & & & & \mathrm{HOT}^{o}_W(S') \arrow[lld, "\sigma'"'] \arrow[dd, "f^{o}_!"] \\
& & \mathrm{HOT}^{\mathrm{lc}\,!}_W(S') \arrow[dd, "f^{\mathrm{lc}\,!}_!"] & & \\
\mathrm{HOT}_W(S) \arrow[rrd, "\rho"] & & & & \mathrm{HOT}^{o}_W(S). \arrow[lld, "\sigma"'] \\
& & \mathrm{HOT}^{\mathrm{lc}\,!}_W(S) & &
\end{tikzcd}
\]

\printedpage{126}
% head: A. Grothendieck : Dérivateurs, Ch. VI, Édition des Universités de Montpellier II et Paris VII
% foot: Version 9 novembre 2001

We do not have a functor $f_!^{\mathrm{lc}}$ at the level of the $\mathrm{HOT}_W^{\mathrm{lc}}$, unless we make an assumption, such as \underline{W(4b)}, which ensures that $\mathrm{HOT}_W^{\mathrm{lc}\,!} \approx \mathrm{HOT}_W^{\mathrm{lc}}$.

The main theorem, that the functor $f_!$ is left adjoint to $f^*$, is valid \textit{under the sole assumption} W$(1,2,3)$ (\footnote{i.e. with these new notations Loc$(1,2,3\ \text{bis})$ (cf. XVI).}). Thus dually for $f_!^{o}$, $f_{o}^*$. But since we do not have a functor $f_{\mathrm{lc}\,!}^*$, the question of adjunction does not arise for general $f$, and without an assumption W(4b) on $W$. It only arises, apparently, for $f$ a trivial $W$-fibration (cf. end of \cercle{6}, p. 198), and in this case the adjunction should not present any difficulties. But as long as I do not have a very clear feeling about the meaning of the $\mathrm{HOT}_W^{\mathrm{lc}\,!}$, this is not a result to get much of a hard-on over \texttt{[sic]}.

The adjunction of $f_!$, $f^*$ for $\mathrm{HOT}_W$ already allows us to show that the prederivator $\mathrm{HOT}_W$
\msp{page 204}
satisfies the conditions Der 1, Der $4'$, Der $5'$ for derivators. Moreover, we have what we can call the property Der 3b for $\mathrm{HOT}_W$ as a $W$-derivator, i.e.

\begin{quote}
\textit{If $f\in W$, then $f$ is a} $\mathrm{HOT}_W$\textit{-equivalence} (\footnote{We prove this by interpreting the notion of a $\mathrm{HOT}_W$-equivalence in terms of $\mathrm{HOT}_W$\textit{-homology.} This comes down in the end to this: If $f : S' \longrightarrow S$ is in $W$, then $\forall\, T$, $S'\times T \longrightarrow S\times T$ too!}) (assuming only the axioms W$(1,2,3)$ on $W$). More precisely, $f$ is a $\mathrm{HOT}_W$-equivalence if and only if $\mathrm{hot}_W(f)$ is \texttt{[an]} isomorphism, i.e. $f\in \overline{W}$, the strong saturation of $W$.
\end{quote}

If $W$ \textit{satisfies} W(4b)\textit{,} then we also have a pair of adjoint functors
\[
\begin{tikzcd}
\mathrm{HOT}_W^{\mathrm{lc}}(S') \arrow[r, shift left, "f_!^{\mathrm{lc}}"] & \mathrm{HOT}_W^{\mathrm{lc}}(S) \arrow[l, shift left, "f^*_{\mathrm{lc}}"]
\end{tikzcd}
\]
for every arrow $S'\longrightarrow S$ in Cat. We also see in this case that % [sic] pout -> pour, misprint

\begin{quote}
% printed with a vertical rule at the left of the block
\begin{tabular}{|p{0.9\linewidth}}
if $f : S' \longrightarrow S$ is in $W$, then $f^*_{\mathrm{lc}}$ and $f_!^{\mathrm{lc}}$ between $\mathrm{HOT}_W^{\mathrm{lc}}(S)$ and $\mathrm{HOT}_W^{\mathrm{lc}}(S')$ are equivalences quasi-inverse to one another (assuming W(4b)),
\end{tabular}
\end{quote}

which is axiom Der 3b$_2$ ``$W$-derivators version''. We cannot expect any better, on the Der 3 side!

\textbf{Omission:} Let $f : S' \longrightarrow S$, we then have $\boxed{f_!f^*} : \mathrm{HOT}_W(S) \longrightarrow \mathrm{HOT}_W(S)$ is equal to $\mathrm{mult}_{S'/S}$ (cf. p. 200). It is therefore equal to $\mathrm{mult}_\xi$, where $\xi = \mathrm{hot}_{S,W}(S')$, if we assume either $S'$ $W$-smooth over $S$, or $S'$ $W$-fibered over $S$ (cf. $(**)$, p. 201).

\subsection*{\protect\cercle{9} Base change theorem for $f_!$ in $\mathrm{HOT}_W$.}

If we have
\printedpage{127}
% head: A. Grothendieck : Dérivateurs, Ch. VI, Édition des Universités de Montpellier II et Paris VII
% foot: Version 9 novembre 2001
\[
\begin{tikzcd}
X \arrow[d, "f"'] & X' \arrow[l, "q"'] \arrow[d, "f'"] \\
S & S', \arrow[l, "g"']
\end{tikzcd}
\]
with $f$ $W$-smooth \textit{or} $g$ $W$-proper \texttt{[the continuation is missing]}